% Generated by code/assemble.py; edit part_B_audit.tex.
\subsection{The power completion and the original recursion}

\begin{result}{E.101}{Proved}{Finite support; SECTIONS\_OVER\_Q A.32 and GT\_TOWER\_CHECK C.1.1 confirmed}
Put $\mathcal K=\varprojlim_n\mathbb Q^\times/(\mathbb Q^\times)^n$,
with discrete power quotients and their inverse-limit topology. We have a
continuous group isomorphism, with continuous inverse,
\[
\mathcal K\simeq\{\pm1\}\times
\left\{(a_q)_q\in\prod_q\widehat{\mathbb Z}:
  \text{for every }n,\ a_q=0\pmod n\text{ for almost all }q\right\}.
\]
The topology on the right is defined by simultaneous reduction of all
coordinates modulo $n$, with the resulting direct sum given the discrete
topology. If $(A,B)$ is ordinary at every finite prime, every valuation
$v_q(A),v_q(B)$ is an ordinary integer. If any one of $A,B,A/B$ has finite
valuation support, there is a unique $t\in\mathbb Q\setminus\{0,1\}$ with
$(A,B)=(\kappa(t),\kappa(1-t))$; all its ordinary local points equal $t$.
No local section theorem is used.
\end{result}
\begin{proof}
Unique factorization identifies $\mathbb Q^\times$ with
$\{\pm1\}\times\bigoplus_q\mathbb Z$. At each power quotient the valuation
coordinates have finite support. A compatible collection of these finite
support vectors gives exactly the displayed restricted family, and vice
versa. Both constructions commute with the defining quotient maps, proving
the assertion about topology as well as the group isomorphism.

The local multiplicative Kummer map
$\mathbb Q_p^\times\to\varprojlim_n\mathbb Q_p^\times/(\mathbb Q_p^\times)^n$
is injective. Indeed its valuation has to be divisible by every $n$, and
is zero. The unit group is a product of a finite group and a group
isomorphic to $\mathbb Z_p$; at two one first uses $1+4\mathbb Z_2$.
The intersection of all power subgroups in each factor is trivial.
Ordinarity therefore gives $v_q(A)=v_q(x_q)\in\mathbb Z$, and similarly
for $B$.

If $A$ has finite support, its sign and integer valuations give an actual
rational number $t$, and the displayed isomorphism gives $A=\kappa(t)$.
Local injectivity gives $x_p=t$ for every $p$, in particular $t\ne0,1$.
The quotient $B/\kappa(1-t)$ is locally trivial everywhere. Its local
valuations annihilate every prime coordinate; its remaining possible sign
is detected at $p=3$, since $-1$ is not a square in $\mathbb Q_3$.
Thus $B=\kappa(1-t)$. Interchanging the coordinates handles finite support
of $B$. If $A/B$ has finite support, it equals $\kappa(u)$ for
$u\in\mathbb Q^\times$, and local injectivity gives
$x_p/(1-x_p)=u$. The value $u=-1$ is impossible, and hence
$x_p=u/(1+u)$ for every $p$. The same local-triviality argument recovers
both global coordinates. Uniqueness follows from any local Kummer map.
\end{proof}

% The following two statements belong in Part B when assembling the paper.

\begin{result}{E.102}{Proved}{SECTIONS\_OVER\_Q A.44 confirmed: ordinary infinite-support pairs}
There are pairwise distinct primes $q_j,r_j>3$, for $j\ge0$, such that
\[
A=3\prod_{j\ge0}q_j^{(j+4)!},\qquad
B=-2\prod_{j\ge0}r_j^{(j+4)!}
\]
converge in $\mathcal K$. At every finite prime $p$ they are the Kummer
classes of nonzero $x_p,y_p\in\mathbb Z_p$ satisfying $x_p+y_p=1$.
Their nonzero valuation vectors are $(1,0)$ at three, $(0,1)$ at two,
$((j+4)!,0)$ at $q_j$, and $(0,(j+4)!)$ at $r_j$. The first prime
$q_0$ can take infinitely many values. The result concerns the abelian
completed pair only; it does not assert a section of any nonabelian
quotient. No local section theorem is used.
\end{result}
\begin{proof}
We retain the explicit factorials of A.44, and independently check the
field intersection and the induction on which its assertion depends.

First let $24\mid h$, let $c=3w^8$ or $c=-2w^8$, with
$w\in\mathbb Z_{>0}$ and $(w,6)=1$, and put
$K=\mathbb Q(\mu_h)$ and $L=K(c^{1/h})$. We claim that the maximal
abelian subextension of $L/\mathbb Q$ is exactly $K$.
The element $c$ is a square in $K$, because $K$ contains both
$\sqrt3$ and $\sqrt{-2}$. For an odd prime $e\mid h$, an $e$th root
of $c$ generates a nonnormal real field of degree $e$: irreducibility
follows from the valuation one at three or two, and the other roots are
nonreal. Such a field cannot lie in a cyclotomic field. A fourth root
of $c$ in $K$ would similarly put $3^{1/4}$ or $2^{1/4}$ in $K$ after
division by the indicated rational square and, for $-2$, a power of
$\zeta_8$. This is again impossible. It follows that the Kummer class
of $c$ has order $h/2$. For detail, a proper prime divisor $e$ of its
possible order would give $c=z^{2e}\eta$ with
$\eta\in\mu_{h/(2e)}$; this root of unity has a $2e$th root in
$\mu_h$, contradicting the exclusions just proved.

Consequently $N=\operatorname{Gal}(L/K)$ is cyclic of order $h/2$.
Its conjugation action is cyclotomic. The subgroup $[N,G]$, where
$G=\operatorname{Gal}(L/\mathbb Q)$, is $N^2$: every cyclotomic
multiplier is odd, and the multiplier $-1$ already generates the ideal
$2$ in $\mathbb Z/(h/2)$. Thus
$N^2\subseteq[G,G]\subseteq N$. The field fixed by $N^2$ is
$K(c^{1/4})$, which contains the nonabelian splitting field of
$X^4-3$ or $X^4-2$. Hence $[G,G]$ cannot be $N^2$, and must be $N$.
This proves the claim.

If $h\mid M$ and $u\in(\mathbb Z/M)^\times$ satisfies
$u=1\pmod h$, the intersection
$L\cap\mathbb Q(\mu_M)$ is therefore $K$.
The two automorphisms $1$ on $L$ and $u$ on $\mathbb Q(\mu_M)$ agree
on that intersection, so define an automorphism of their compositum.
Chebotarev's theorem \citep[Theorem 8.31, p.~147]{MilneANT2020} supplies
infinitely many primes in the class $u\pmod M$ splitting completely
in $L$. We may exclude any finite set. This is the precise simultaneous
prime-selection statement used below.

For the local existence bound, put $d=(m,p-1)$ and assume $p\nmid mab$.
Character orthogonality on $\mathbb F_p^\times$ gives at least
\[
p-2-3(d-1)-(d-1)(d-2)\sqrt p
\]
nonzero solutions of $aX^m+bY^m=1$. Indeed the trivial character pair
contributes $p-2$, the $3(d-1)$ terms with one trivial character or
trivial product have absolute value one, and each remaining Jacobi sum
has absolute value $\sqrt p$. The latter identity follows by expanding
the two Gauss sums and changing variables $(x,y)=(zt,z(1-t))$; the
absolute value of each nontrivial Gauss sum is $\sqrt p$ by
orthogonality. If $p>(m+1)^4$, the displayed lower bound is positive:
with $T=\sqrt p$, the lower polynomial
$T^2-(m-1)(m-2)T-3m+1$ is increasing beyond $(m+1)^2$ and equals
$m^2(5m+9)>0$ there. Such a point has a unit derivative, and lifts.

We also use the following sharpened Newton estimate. If
$f\in\mathbb Z_p[T]$ and $v_p(f(1))>2e$, where
$e=v_p(f'(1))$, Newton iteration gives a root $z$ with
$v_p(z-1)\ge v_p(f(1))-e$. The first Newton increment has that
valuation, the Taylor remainder at least twice that valuation, and the
derivative keeps valuation $e$; all subsequent increments have larger
valuation. This proves both convergence and the stated bound, including
$p=2$ (compare \citealp[Theorem 7.32 and the sharper bound on p.~121]{MilneANT2020}).

Set $n_j=(j+4)!$, $B(n)=(n+1)^4$, and
$H_j(p)=2v_p(n_{j+1})+j+3$. We inductively construct
\[
a_j=3u_j^8,\quad b_j=-2v_j^8,\quad (u_jv_j,6)=1,
\quad (a_j,b_j)=1,
\]
such that $a_jX^{n_j}+b_jY^{n_j}=1$ has a unit-coordinate point at
every finite prime and
$v_p(a_j+b_j-1)\ge H_j(p)$ for every $p\le B(n_j)$.
The initial coefficients $(3,-2)$ satisfy all these conditions.

Suppose they hold at stage $j$, and write $n=n_j$, $m=n_{j+1}$,
$a=a_j$, $b=b_j$. Let $S$ contain all primes at most $B(m)$ and all
prime divisors of $ab$. Put $L_p=H_{j+1}(p)$ and
$M=\prod_{p\in S}p^{L_p}$; thus $m\mid M$ and
$L_p>2v_p(m)$. Choose old unit-coordinate solutions
$(\xi_p,\eta_p)$ at every $p\in S$. At $p\mid m$, choose them
congruent to $(1,1)$ modulo $p^{v_p(m)}$: the old error bound applies
there, and Newton in a variable with unit coefficient gives precisely
this congruence. At the other old primes any unit solution suffices.

The CRT unit class defined by the $\xi_p$ is one modulo $m$.
The prime-selection statement, applied to $\mathbb Q(\mu_n,b^{1/n})$,
gives a new prime $q$ in this class, larger than every prime of $S$
and than $n_{j+2}$. Since $b$ is an $n$th power modulo $q$,
simple-root lifting gives a unit $z_q$ with
$bz_q^n=1-aq^n$. Impose next the congruences
$r=\eta_p\pmod{p^{L_p}}$ for $p\in S$ and
$r=z_q\pmod q$. Their CRT class is still one modulo $m$.
Apply the same prime-selection statement with $h=m$ and
$c=aq^n$, which still has the required eighth-power form, to obtain a
prime $r>q$ splitting in $\mathbb Q(\mu_m,(aq^n)^{1/m})$.
Define $a_{j+1}=aq^n$, $b_{j+1}=br^n$.

At every $p\in S$ their sum is one modulo $p^{L_p}$, so Newton gives
a unit-coordinate solution of the new exponent-$m$ equation and the
next error bound. At $q$ the new equation holds at $(1,1)$ modulo $q$,
and its $Y$ derivative is a unit. At $r$ the coefficient $aq^n$ is an
$m$th power modulo $r$; setting $Y=1$ gives an equation for $X$ with a
simple nonzero root. All other primes are larger than $B(m)$ and divide
neither coefficient, so the character-sum bound applies. These are all
finite primes, completing the induction. Both prime choices can always
avoid all previous primes, and already at its first use the choice of
$q$ ranges over an infinite set.

For each fixed power modulus, the factorial exponents eventually
vanish modulo it. The coefficient products therefore converge in
$\mathcal K$ and have exactly the asserted valuation vectors. At a
fixed finite prime $p$, the valuations become constant. Subsequent
ratios are $n_j$th powers of local units; the factorials eventually
annihilate every finite quotient of $\mathbb Z_p^\times$. The two
sequences thus converge in $\mathbb Q_p^\times$ to nonzero elements
with those constant valuations. The error valuation tends to infinity,
so their sum is one. Openness of local power subgroups, also a
consequence of Newton's estimate, identifies these ordinary limits
with the localizations of the completed products. The real signs are
$+,-$, represented by the ordinary pair $(3,-2)$. Infinite support
excludes a rational pair by E.101. Taking at each stage the least
eligible residue choices and prime choices specifies an actual
terminating recursion; no numerical execution of its huge first
prime search is being asserted.
\end{proof}


\begin{result}{E.103}{Known}{SECTIONS\_OVER\_Q A.42 confirmed from rational-point finiteness}
The rational unit-equation graph
$\mathcal R=\{(\kappa(t),\kappa(1-t)):t\in\mathbb Q\setminus\{0,1\}\}$
is closed and discrete in $\mathcal K^2$. Every Cauchy net in it is
eventually constant. Thus a nonrational completed pair is separated
from all rational pairs in some simultaneous finite power quotient.
\end{result}
\begin{proof}
In a fixed fourth-power fibre choose representatives $a,b$.
Every rational pair in that fibre comes from a rational point of the
smooth projective quartic $aX^4+bY^4=Z^4$. Its genus is three, so
\citet[Satz~7, p.~365]{Faltings1983} makes this fibre finite. Each fourth-power
fibre is open and closed, and the ambient completion is Hausdorff.
Removing finitely many points gives both closedness and discreteness
of the graph. A Cauchy net eventually lies in one such finite fibre;
a common power quotient separates its finitely many points, so the
net is eventually constant. Basic open neighbourhoods are simultaneous
power-quotient fibres, proving the last assertion.
\end{proof}

